\chapter{Conclusiones}
\label{chap:conclusiones}

Tras el diseño, la simulación y la implementación física del sistema de procesamiento de señales, se concluye que los
objetivos planteados fueron alcanzados satisfactoriamente. El análisis de los resultados permite extraer las siguientes
observaciones finales:

\begin{itemize}
    \item La arquitectura modular basada en etapas secuenciales demostró ser robusta, permitiendo la corrección puntual
    de parámetros en las etapas 3 y 6 sin comprometer el funcionamiento global del sistema.
    \item El uso de operacionales con entrada JFET (TL084) fue clave para minimizar los errores por corrientes de
    polarización en los integradores, garantizando una señal triangular limpia y con pendientes constantes.
    \item La implementación física puso de manifiesto que el uso de una protoboard resultó problemático por los falsos
    contactos. Esta limitación técnica generó dificultades para estabilizar las mediciones y reafirma la importancia de
    considerar montajes más permanentes (como PCB) para circuitos con múltiples etapas de ganancia y realimentación.
    \item Se destaca que priorizar el uso de resistencias bajo el estándar E24 fue favorecedor al momento de tener que
    implementar todo. La facilidad para encontrar valores comerciales que se ajustaran a los cálculos de diseño redujo
    el tiempo de montaje y facilitó el ajuste de las ganancias finales.
    \item La discrepancia observada entre la simulación y la práctica subraya la necesidad de incluir márgenes de ajuste
    (como el uso de trimpots) o de realizar cálculos redundantes rápidas basadas en componentes reales para absorber las
    tolerancias de fabricación.
\end{itemize}

En definitiva, el trabajo práctico permitió consolidar los conceptos de diseño con amplificadores operacionales en
configuraciones lineales y no lineales, enfrentando con éxito los desafíos inherentes a la transición del modelo teórico
al prototipo físico.
